%Dokumentklasse
\documentclass[a4paper,12pt]{scrreprt}
\usepackage[left= 2.5cm,right = 2cm, bottom = 4 cm]{geometry}
 \usepackage[onehalfspacing]{setspace}
% ============= Packages =============

% Dokumentinformationen
\usepackage[
	pdftitle={Testautomatisierung},
	pdfsubject={},
	pdfauthor={Matthias Karl},
	pdfkeywords={}
	pdftex=true, 
	colorlinks=true,
 	breaklinks=true,
	citecolor=black,
	linkcolor=black,	
	menucolor=black,	
	urlcolor=black
]{hyperref}


% Standard Packages
\usepackage[utf8]{inputenc}
\usepackage[ngerman]{babel}
\usepackage[T1]{fontenc}
\usepackage{graphicx}
\graphicspath{{img/}}
\usepackage{fancyhdr}
\usepackage{lmodern}
\usepackage{color}
\usepackage{transparent}

% zusätzliche Schriftzeichen der American Mathematical Society
\usepackage{amsfonts}
\usepackage{amsmath}

% nicht einrücken nach Absatz
\setlength{\parindent}{0pt}


% ============= Kopf- und Fußzeile =============
\pagestyle{fancy}
%
\lhead{}
\chead{}
\rhead{\slshape \leftmark}
%%
\lfoot{}
\cfoot{}
\rfoot{\thepage}
%%
\renewcommand{\headrulewidth}{0.4pt}
\renewcommand{\footrulewidth}{0pt}

% ============= Package Einstellungen & Sonstiges ============= 

%Besondere Trennungen
\hyphenation{De-zi-mal-tren-nung St-rei-fen-licht-scan-nern}

%römische Aufzählungen mit \RM{Zahl}
\newcommand{\RM}[1]{\MakeUppercase{\romannumeral #1}}


% ============= Dokumentbeginn =============

\begin{document}

\pagestyle{empty}
\begin{center}
\begin{tabular}{p{\textwidth}}


\begin{center}
\includegraphics[scale=0.17]{img/logos.jpg}
\end{center}


\\

\begin{center}
\LARGE{\textsc{
Testautomatisierung \\  \large{- mit Selenium -}
}}
\end{center}

\\


\begin{center}
\large{Fakultät für Informatik und Mathematik \\
der Hochschule München \\}
\end{center}

\\

\begin{center}
\textbf{\Large{Seminararbeit}}
\end{center}



\begin{center}
vorgelegt von
\end{center}

\begin{center}
\large{\textbf{Matthias Karl}} \\
\large{Matrikel-Nr: 03280712} \\
\end{center}

\begin{center}
\large{im Dezember 2014}
\end{center}

\\

\\

\begin{center}
\begin{tabular}{lll}
\textbf{Prüferin:} & & Prof. Dr. Ulrike Hammerschall\\
\end{tabular}
\end{center}

\end{tabular}
\end{center}

% \part im Inhaltsverzeichnis nicht nummerieren
\makeatletter
\let\partbackup\l@part
\renewcommand*\l@part[2]{\partbackup{#1}{}}

%Seitennummerierung neu beginnen, Zahlen [arabic], röm.Zahlen [roman,Roman], Buchstaben [alph,Alph]
\pagenumbering{Roman}

\newpage
\addsec{Eidesstattliche Erklärung}
\label{erklaerung}

Hiermit versichere ich, die vorliegende Studienarbeit selbstständig und nur unter Verwendung der von mir angegebenen Quellen und Hilfsmittel verfasst zu haben. Sowohl inhaltlich als auch wörtlich entnommene Inhalte wurden als solche kenntlich gemacht. Die Arbeit hat in dieser oder vergleichbarer Form noch keinem anderem Prüfungsgremium vorgelegen. \\
\\[1.5cm]
Datum:	\hrulefill\enspace Unterschrift: \hrulefill
\\[3.5cm]
 



\addsec{Zusammenfassung / Abstract}
\label{sec:zusammenfassung}
Diese Arbeit befasst sich mit der Automatisierung von Testfällen. Die Arbeit gibt einen Überblick über die verschiedenen Bereiche und Möglichkeiten bei der Testautomatisierung und befasst sich im nähren mit dem Automatisierungstool Selenium.
Zu Beginn werden zunächst die allgemeinen Grundlagen des Testens geklärt um
dann die verschiedenen Bereiche und Möglichkeiten der Testautomatisierung aufzuzeigen. Auf den Bereich der automatisierten UI-Tests wird im Kontext von Selenium zum Schluss näher eingegangen.



\newpage
\pagestyle{fancy}
%Inhaltsverzeichnis
\tableofcontents

\newpage
%Seitennummerierung neu beginnen, Zahlen [arabic], röm.Zahlen [roman,Roman], Buchstaben [alph,Alph]
\pagenumbering{arabic}
% pagestyle für gesamtes Dokument aktivieren
\pagestyle{fancy}

\newpage
\chapter{Einleitung}
\label{sec:einleitung}

Software hat in der heutigen Zeit eine hohe Verbreitung gefunden. Softwaresysteme werden immer wichtiger. Sowohl für Unternehmen als auch für jeden einzelnen persönlich. Die Anforderungen an moderne Software steigen ständig. Die Systeme werden immer größer und komplexer. Das hat zur Folge, dass auch die Anforderungen an die Qualität der Software immer weiter Steigen und wichtiger werden.
Fehler in Software verursachen immer wieder einen hohen finanziellen Schaden und können im schlimmsten Fall sogar Menschenleben kosten. Diese Probleme werden immer gravierender wenn mit der Wichtigkeit, Komplexität und Größe von Software nicht auch gleichzeitig die Qualität steigt. \cite{burnstein_practical_2003} \newline
Softwaretests sind ein weit verbreitetes Mittel um die Qualität einer Software zu überprüfen und sicherzustellen.
Komplexere und größere Software bedeutet daher auch gleichzeitig einen steigenden Testaufwand.\newline
Qualitätssicherungsmaßnahmen wie beispielsweise das Testen machen jetzt schon einen Großteil der Kosten in Softwareprojekten aus. Studien haben gezeigt, dass das Testen für 50\% und mehr der gesamten Projektkosten verantwortlich ist. \cite{ramler_economic_2006} 
Es ist daher nicht verwunderlich, dass es immer wieder versuche gibt diese Kosten zu reduzieren ohne dabei den angestrebten Qualitätsstandart zu reduzieren.
Einer der Wege die dafür vorgeschlagen wurden ist das automatisieren von Testfällen 
\cite{harrold_testing:_2000}
Im laufe der Jahre hat die Teatautomatisierung immer mehr an Bedeutung gewonnen. Heute ist sie bereits fester Bestandteil von Bereichen wie Continous Delivery und Continous Integration.

Diese Arbeit befasst sich daher mit der Automatisierung von Testfällen. Die Arbeit soll einen Überblick über die verschiedenen Bereiche der Testautomatisierung geben um dann den Bereich der Automatisierten GUI-Tests näher zu beleuchten.
Kapitel 1 ...






\newpage
\chapter{Grundlagen}
\label{sec:grundlagen}


\section{Softwarequalität}
\label{sec:softwarequalität}

Nach der Norm ISO-25000:2014 4.33 bezeichnet der Begriff der Softwarequalität die Fähigkeit einer Software die expliziten und impliziten Bedürftnisse von Benutzern, unter den Bedingungen unter sie benutzt wird, zu befriedigen.  \cite {international_organization_for_standardization_iso_so/iec_2014}
Softwarequalität hat nach dieser Definition einen subjektiven Charakter.
Dieser subjektive Charakter macht den Begriff in der Praxis schwer zu greifen und damit nicht direkt anwendbar. Aus diesem Grund existieren sogenannte Qualitätsmodelle, die den Begriff der Softwarequalität messbar und damit auch überprüfbar machen sollen.
Ein solches Qualitätsmodell wird zum Beispiel in der ISO-Norm 9126 vorgestellt.
Es werden verschiedene Qualitätsmerkmale definiert die zur Beurteilung der Gesamtqualität eines Softwareprodukts dienen.
Hierunter fallen die Merkmale:

\begin{itemize}
	  \itemsep0pt
      \item Funktionalität
      \item Zuverlässigkeit
      \item Benutzbarkeit
      \item Effizienz
      \item Änderbarkeit
      \item Übertragbarkeit          
\end{itemize}

Es existieren verschiedene Methoden um sicherzustellen, dass Software bezogen auf die Qualitätsmerkmale gewissen Anforderungen genügt.
Eine Teil der Methoden geht dabei davon aus, dass ein qualitativ hochwertiger Prozess der Produkterstellung die Entstehung von qualitativ hochwertigen Produkten begünstigt.
Das Augenmerk wird hierbei also auf die Prozessqualität gelegt.
Allgemein fasst man diese Gruppe unter dem Begriff des prozessorientiertes Qualitätsmanagement zusammen. Die klassischen Vorgehensmodelle der Softwarentwicklung werden z.B. hier eingeordnet.
Worauf sich diese Arbeit jedoch konzentrieren möchte sind die Methoden des produktorientierten Qualitätsmanagement.
Hierbei wird das Softwareprodukt direkt bezüglich der Qualitätsmerkmale überprüft.
Darunter fallen beispielsweise Softwaretests.


\section{Softwaretest}
\label{sec:softwaretest}

Das produktorientierte Qualitätsmanagement unterteilt sich weiter in die Bereiche des konstruktiven und analytischen Qualitätsmanagement.
Unter dem konstruktiven Qualitätsmanagement versteht man in diesem Fall den Einsatz von z.B. Methoden, Werkzeugen oder Standards die dafür sorgen, dass ein (Zwischen-)Produkt bestimmte Forderungen erfüllt.
Was man im Allgemeinen aber unter einem Softwaretest versteht ist im Bereich der prüfenden Verfahren des analytischen Qualitätsmanagement angesiedelt.
Unter analytischen Qualitätsmanagement versteht man hier den Einsatz von analysierenden bzw. prüfenden Verfahren, die Aussagen über die Qualität eines (Zwischen-)Produkts machen. \newline
Die Norm ISO/IEC/IEEE 24765:2010 3.280 definiert das testen von Software als \glqq the dynamic verification of the behavior of a program on a finite set of test cases, suitably selected from the usually infinite executions domain, against the expected behavior.\grqq\
Aufgabe eines Softwaretests ist es dabei nicht einen Fehler im Code zu Lokalisieren und zu beheben. TODO:CITE
Das Lokalisieren und Beheben des Defekts ist Aufgabe des Softwareentwicklers und wird auch als Debugging (Fehlerbereinigung, Fehlerkorrektur) bezeichnet.
Während Debugging das Ziel hat, Defekte bzw. Fehlerzustände zu beheben, ist es Aufgabe des Tests, Fehlerwirkungen (die auf Defekte hinweisen) gezielt und systematisch aufzudecken. \cite{spillner_basiswissen_2003}
Dabei dienen definierte Anforderungen als Prüfreferenz, mittels derer ggf. vorhandene Fehler aufgedeckt werden.
\glqq Das Testen von Software dient durch die Identifizierung von Defekten und deren anschließenden Beseitigung durch das Debugging zur Steigerung der Softwarequalität\grqq\ \cite{spillner_basiswissen_2003}
Als möglicher Rahmen für die Anforderungen können z.B. die in \ref{sec:softwarequalität} bereits beschriebenen Qualitätsmerkmale dienen.


\section{Testprozess}
\label{sec:testprozess}

Der Begriff des Softwaretests wie er in \ref{sec:softwaretest} beschrieben ist erfordert eine Einordnung in einen größeren Zusammenhang. Ein Softwaretest steht in der Regel nicht für sich alleine, sondern ist Teil eines größeren Prozesses der den Softwaretest in seinem gesamten Lebenszyklus begleitet.
Durch den Testprozess wird die Aufgabe des Testens in kleinere Arbeitsschritte gegliedert.
Die Phasen die man dabei unterscheidet sind:

\begin{itemize}
	  \itemsep0pt
      \item Testplanung und Steuerung
      \item Testanalyse und -design
      \item Testrealisierung und -durchführung
      \item Testauswertung und Berischt sowie Abschluss der Testaktivitäten       
\end{itemize}

Obgleich die Aufgaben in sequenzieller Reihenfolge im Testprozess angegeben sind, können sie sich überschneiden und teilweise auch gleichzeitig durchgeführt werden. Diese Teilaufgaben bilden einen Fundamentalen Testprozess und werden im Folgenden näher beschrieben. \cite[Seite 19]{spillner_basiswissen_2003}
\newpage
\chapter{Testautomatisierung}
\label{sec:testautomatisierung}

Nach Seidl et al. versteht man unter dem Begriff der Testautomatisierung \glqq die Durchführung von ansonsten manuellen Testtätigkeiten durch Automaten.\grqq\ \cite[Seite 7]{seidl_basiswissen_2012}
Das Spektrum der Testautomatisierung umfasst demnach alle Tätigkeiten die dazu dienen die Qualität einer Software zu Überprüfen. Darunter fallen alle Aufgaben die im Testprozess in den einzelnen Phasen des Softwareentwicklung anfallen.
Testautomatisierung ist also nicht nur auf die automatisierte Testdurchführung beschränkt, sondern kann ebenso bei der Testfallerstellung, der Testdatengenerierung, der Testauswertung oder auch der Testumgebungsherstellung und -wiederherstellung eine Rolle spielen.
\glqq Die Grenze der Automatisierung liegen [nach Seidl et al.] darin, dass diese nur die manuellen Tätigkeiten eines Testers übernehmen kann, nicht aber die intellektuelle, kreative und intuitive Dimension dieser Rolle.\grqq\ \cite[Seite 7]{seidl_basiswissen_2012}
Besonders lohnenswert ist eine Automatisierung bei sich wiederholenden Aufgaben, zu denen
besonders die Testdurchführung zählt. Die Automatisierung in diesem Bereich ist bereits weit verbreitet und soll auch den Fokus dieser Arbeit bilden. Die anderen Bereiche der Testautomatisierung sollen nur kurz aufgezeigt werden.


%Verzeichnis aller Bilder
\newpage
\listoffigures

%Literaturverzeichnis
\newpage
\bibliographystyle{alphadin}
\bibliography{Literatur}

\end{document}
